%-------------------------
% Резюме на LaTeX
% Автор : адаптировано для Рахима Баймурзина
% Сборка: tectonic cv-ru.tex   (или)  xelatex cv-ru.tex
%          Нужен движок с поддержкой Unicode (XeTeX/LuaTeX).
%------------------------

\documentclass[letterpaper,11pt]{article}

\usepackage[empty]{fullpage}
\usepackage{titlesec}
\usepackage[usenames,dvipsnames]{color}
\usepackage{enumitem}
\usepackage[hidelinks]{hyperref}
\usepackage{fancyhdr}
\usepackage{comment}
\usepackage{array}
\usepackage{iftex}

% Computer Modern Unicode — тот же рисунок шрифта, но с кириллицей.
\usepackage{fontspec}
\setmainfont{cmunrm.otf}[
  BoldFont          = cmunbx.otf,
  ItalicFont        = cmunti.otf,
  BoldItalicFont    = cmunbi.otf,
  SmallCapsFont     = cmunrm.otf,
  SmallCapsFeatures = {Letters=SmallCaps},
]

\usepackage{polyglossia}
\setmainlanguage{russian}
\setotherlanguage{english}

% Карта глифов для тегированного PDF: только pdfTeX, под XeTeX/LuaTeX пропускается.
\ifPDFTeX
  \input{glyphtounicode}
  \pdfgentounicode=1
\fi

\pagestyle{fancy}
\fancyhf{}
\fancyfoot{}
\renewcommand{\headrulewidth}{0pt}
\renewcommand{\footrulewidth}{0pt}

\addtolength{\oddsidemargin}{-0.5in}
\addtolength{\evensidemargin}{-0.5in}
\addtolength{\textwidth}{1in}
\addtolength{\topmargin}{-.6in}
\addtolength{\textheight}{1.3in}

\urlstyle{same}
\raggedbottom
\raggedright
\setlength{\tabcolsep}{0in}

\titleformat{\section}{
  \vspace{-4pt}\scshape\raggedright\large
}{}{0em}{}[\color{black}\titlerule \vspace{-5pt}]

%-------------------------
% Макросы резюме
%-------------------------

\newcommand{\resumeItem}[1]{
  \item\small{
    {#1 \vspace{-2pt}}
  }
}

\newcommand{\resumeSubheading}[4]{
  \vspace{-2pt}\item
    \begin{tabular*}{0.97\textwidth}[t]{l@{\extracolsep{\fill}}r}
      \textbf{#1} & #2 \\
      \textit{\small#3} & \textit{\small #4} \\
    \end{tabular*}\vspace{-7pt}
}

\newcommand{\resumeProjectHeading}[2]{
    \item
    \begin{tabular*}{0.97\textwidth}{l@{\extracolsep{\fill}}r}
      \small#1 & #2 \\
    \end{tabular*}\vspace{-7pt}
}

\renewcommand\labelitemii{$\vcenter{\hbox{\tiny$\bullet$}}$}


\newcommand{\resumeSubHeadingListStart}{\begin{itemize}[leftmargin=0.15in, label={}]}
\newcommand{\resumeSubHeadingListEnd}{\end{itemize}}
\newcommand{\resumeItemListStart}{\begin{itemize}}
\newcommand{\resumeItemListEnd}{\end{itemize}\vspace{-5pt}}

%-------------------------------------------
%%%%%%  РЕЗЮМЕ НАЧИНАЕТСЯ ЗДЕСЬ  %%%%%%%%%%%%%%%%%%%%%%%%%%%%

\begin{document}

\begin{center}
    \textbf{\Huge \scshape Рахим Баймурзин} \\
    Казахстан $|$
    \href{mailto:baimurzzin@gmail.com}{\underline{baimurzzin@gmail.com}} $|$
    \href{https://www.linkedin.com/in/baimurzzin/}{\underline{linkedin}} $|$
    tg: @baimurzzin
\end{center}

%-----------ОБРАЗОВАНИЕ-----------
\section{Образование}
\resumeSubHeadingListStart
  % \resumeSubheading
  %   {Назарбаев Интеллектуальная школа (НИШ)}{Караганда, Казахстан}
  %   {Аттестат о среднем образовании, GPA 3.75, ЕНТ 138/140, первый в истории школы медалист ISO}{2016 -- 2022}
  \resumeSubheading
    {Назарбаев Университет}{Астана, Казахстан}
    {Магистр прикладной математики}{2026 -- 2028}
  \resumeSubheading
    {Назарбаев Университет}{Астана, Казахстан}
    {Бакалавр математики, GPA 3.21}{2022 -- 2026}
\resumeSubHeadingListEnd

%-----------НАГРАДЫ-----------
\section{Награды и достижения}
\resumeSubHeadingListStart
  \item{\begin{tabular}{@{}>{\raggedright\arraybackslash}p{0.45\textwidth}@{\hspace{0.05\textwidth}}>{\raggedright\arraybackslash}p{0.45\textwidth}@{}}
    \textbf{IMO 2021, бронза}                       & \textbf{Региональный Финалист ICPC (NEF)}\\
    IMC 2026, серебро                               & Республиканская Студенческая МО 2026 (1/50)\\
    MBZUAI K2 Think Хакатон, победитель (1/900)     & KAIST E5 Global Конкурс Стартапов (2/40)
  \end{tabular}}
\resumeSubHeadingListEnd

%-----------ОПЫТ РАБОТЫ-----------
\section{Опыт работы}
\resumeSubHeadingListStart

  \resumeSubheading
    {Desargues AI}{Астана, Казахстан}
    {Сооснователь}{фев 2026 -- сейчас}
    \resumeItemListStart
      \resumeItem{Разрабатываю Верифицируемый Математический Интеллект / Математический Суперинтеллект.}
    \resumeItemListEnd

  \resumeSubheading
    {Higgsfield AI}{Алматы, Казахстан}
    {ML-инженер}{сен 2025 -- фев 2026}
    \resumeItemListStart
      \resumeItem{Собирал датасеты и делал высокочувствительные классификаторы для сгенерированных и реальных изображений в проде.}
    \resumeItemListEnd

  \resumeSubheading
    {Mercor AI Lab}{Удалённо}
    {Эксперт по олимпиадной и научной математике / инженер данных}{2025}
    \resumeItemListStart
      \resumeItem{Собирал и оценивал данные для создания лучшей модел  в рейтинге FrontierMath.}
    \resumeItemListEnd

  \resumeSubheading
    {Математические олимпиады}{Казахстан}
    {Тренер по математике, организатор сборов, жюри, автор задач}{2021 -- сейчас}
    \resumeItemListStart
      \resumeItem{Проводил учебно-тренировочные сборы для участников национальной команды IMO.}
      \resumeItem{Подготовил 4 учеников с базового уровня до уровня международных медалистов.}
      \resumeItem{Жюри и организатор на 20+ региональных, республиканских и международных олимпиадах.}
    \resumeItemListEnd

\resumeSubHeadingListEnd

%-----------АКАДЕМИЧЕСКИЙ ОПЫТ-----------
\section{Академический опыт}
\resumeSubHeadingListStart

  \resumeSubheading
    {Назарбаев Университет}{Астана, Казахстан}
    {Ассистент преподавателя по микроэкономике}{осень 2024}
    \resumeItemListStart
      \resumeItem{Вёл семинары и помогал студентам готовиться к экзаменам.}
    \resumeItemListEnd

  \vspace{-18pt}

  \resumeSubheading
    {}{}
    {Научный ассистент по математике}{осень 2025}
    \resumeItemListStart
      \resumeItem{Работал с профессорами Рустемом Тахановым и Женисбеком Асылбековым над критериями соответствия распределений. Создавал тестовые распределения для оценки критериев.}
    \resumeItemListEnd

  \vspace{-18pt}

  \resumeSubheading
    {}{}
    {Исследователь}{весна 2026}
    \resumeItemListStart
      \resumeItem{Работал под руководством профессора Адильбека Каиржана над дипломной работой по дисперсионным уравнениям в частных производных.}
    \resumeItemListEnd

  \vspace{-18pt}

  \resumeSubheading
    {}{}
    {Волонтёр 15-го конгресса ISAAC}{август 2025}
    \resumeItemListStart
      \resumeItem{Сопровождал иностранных гостей, вёл интервью, организовывал и посещал семинары.}
    \resumeItemListEnd

  \resumeSubheading
    {Imperial College London}{Лондон, Великобритания}
    {Эксперт по формализации математики (LEAN)}{июль 2026}
    \resumeItemListStart
      \resumeItem{Участвовал в практикуме по формализации великой теоремы Ферма под руководством Кевина Баззарда. Формализовал теорию представлений.}
    \resumeItemListEnd

\resumeSubHeadingListEnd

% %-----------ВОЛОНТЁРСКИЙ ОПЫТ-----------
% \section{Волонтёрский опыт}
% \resumeSubHeadingListStart
%
%   \resumeSubheading
%     {NU SIAM Student Chapter}{Астана, Казахстан}
%     {Участник и представитель математиков}{2022 -- 2026}
%     \resumeItemListStart
%       \resumeItem{Организовывал математические доклады, соревнования (например, Integration Bee) и представлял студентов-математиков на факультетских обсуждениях.}
%       \resumeItem{Создал и вёл академический Telegram-канал с 1000+ подписчиков.}
%     \resumeItemListEnd
%
% \resumeSubHeadingListEnd

%-----------КУРСЫ-----------
% \section{Профильные курсы}
% \resumeSubHeadingListStart
%   \item{\begin{tabular}{@{}p{0.48\textwidth} p{0.48\textwidth}@{}}
%     MATH 465: дифференциальная геометрия & MATH 481: уравнения в частных производных \\
%     MATH 510: теория меры                & MATH 602: аналитическая теория чисел      \\
%     MATH 682: функциональный анализ      & CSCI 371: алгоритмы II                    \\
%     CSCI 694: глубокое обучение          & ECON 498: общее равновесие                \\
%     CHEM 211: органическая химия         & MINE 402: добыча угля
%   \end{tabular}}
% \resumeSubHeadingListEnd

\begin{comment}
%-----------ПРОЕКТЫ-----------
\section{Проекты}
\resumeSubHeadingListStart

  \resumeProjectHeading
    {\textbf{Детекция рака кожи}}{\textit{осень 2025}}
    \resumeItemListStart
      \resumeItem{Разработал модель на основе CNN для выявления рака кожи по медицинским изображениям.}
      \resumeItem{Исследовал применение глубокого обучения в медицинской диагностике.}
    \resumeItemListEnd

\resumeSubHeadingListEnd
\end{comment}

%-----------НАВЫКИ И УВЛЕЧЕНИЯ-----------
\section{Прочее}
\resumeSubHeadingListStart
  \resumeItem{\textbf{Навыки:} ML, DS, GitHub, C++, Python, AWS, Docker, SQL, BigQuery, Grafana, LaTeX, Excel, Apify.}
  % \resumeItem{\textbf{Увлечения:} спортивное программирование (Codeforces 1600), шахматы (1700 блиц/рапид), фортепиано, настольный теннис.}
\resumeSubHeadingListEnd

\end{document}
